% ===== PRÉAMBULE CANONIQUE — à coller VERBATIM en tête de CHAQUE .tex =====
\documentclass[11pt,a4paper]{article}
\usepackage[utf8]{inputenc}\usepackage[T1]{fontenc}\usepackage{lmodern}
\usepackage[french]{babel}
\usepackage{amsmath,amssymb,mathtools}\usepackage{icomma}
\usepackage[version=4]{mhchem}          % \ce{...} : équations & formules
\usepackage{siunitx}\sisetup{locale=FR} % \qty{}{}, \si{}, \ang{}
\usepackage{chemfig}                    % \chemfig{...} : structures développées
\usepackage{xcolor}\usepackage{colortbl}
\usepackage{tikz}\usepackage{pgfplots}\pgfplotsset{compat=1.18}
\usepackage[most]{tcolorbox}\usepackage{enumitem}\usepackage{array,booktabs}
\usepackage[a4paper,margin=2.2cm]{geometry}
\usetikzlibrary{arrows.meta,calc,angles,quotes,positioning,patterns,backgrounds,shapes.geometric}
\definecolor{primary}{RGB}{222,49,99}\definecolor{primarydark}{RGB}{150,22,55}
\definecolor{defcol}{RGB}{0,114,178}\definecolor{methcol}{RGB}{52,92,148}
\definecolor{excol}{RGB}{0,135,90}\definecolor{attncol}{RGB}{200,40,55}
\tcbset{boxbase/.style={enhanced,breakable,boxrule=0.9pt,arc=2.5pt,
  left=3mm,right=3mm,top=2.3mm,bottom=2.3mm,fonttitle=\bfseries,coltitle=white}}
% --- boîtes TUTORIEL (noms EXACTS, ne pas renommer) ---
\newtcolorbox{cle}[1][]{boxbase,colback=defcol!6,colframe=defcol,
  title={\textsc{À retenir}\ifx\relax#1\relax\relax\else\textmd{ : \itshape #1}\fi}}
\newtcolorbox{methode}[1][]{boxbase,colback=methcol!7,colframe=methcol,sharp corners,
  title={\textsc{Méthode}\ifx\relax#1\relax\relax\else\textmd{ --- \itshape #1}\fi}}
\newtcolorbox{exemplebox}[1][]{boxbase,colback=excol!5,colframe=excol,
  title={\textsc{Exemple}\ifx\relax#1\relax\relax\else\textmd{ : \itshape #1}\fi}}
\newtcolorbox{piege}[1][]{boxbase,colback=attncol!6,colframe=attncol,
  title={\textsc{Piège}\ifx\relax#1\relax\relax\else\textmd{ : \itshape #1}\fi}}
\newtcolorbox{remarque}[1][]{boxbase,colback=black!4,colframe=black!45,
  title={\textsc{Remarque}\ifx\relax#1\relax\relax\else\textmd{ : \itshape #1}\fi}}
% --- boîtes EXERCICE / CORRIGÉ (noms EXACTS, auto-comptées) ---
\newtcolorbox[auto counter]{exo}[1][]{boxbase,colback=primary!4,colframe=primary,
  title={\textsc{Exercice}~\thetcbcounter\ifx\relax#1\relax\relax\else\textmd{ --- \itshape #1}\fi}}
\newtcolorbox[auto counter]{corrige}[1][]{boxbase,colback=excol!4,colframe=excol,
  title={\textsc{Corrigé}~\thetcbcounter\ifx\relax#1\relax\relax\else\textmd{ --- \itshape #1}\fi}}
\newcommand{\R}{\mathbb{R}}\newcommand{\N}{\mathbb{N}}\newcommand{\Z}{\mathbb{Z}}\newcommand{\ds}{\displaystyle}
% (un \usepackage{fancyhdr}+en-tête est possible ; il est ignoré côté web)
% =========================================================================
\begin{document}
\sisetup{per-mode=symbol}
\begin{exo}[la soustraction qui détruit la précision]
Effectue chaque soustraction et donne le résultat avec le bon nombre de décimales. Combien de CS
reste-t-il ?
\begin{enumerate}[label=\alph*)]
  \item $8,11 - 8,0$
  \item $100,0 - 99,7$
  \item $5,432 - 5,41$
\end{enumerate}
\end{exo}

\begin{exo}[calcul mixte : masse molaire puis moles]
On donne $M(\ce{H})=\qty{1,008}{\gram\per\mole}$ et $M(\ce{O})=\qty{16,00}{\gram\per\mole}$. Les indices
d'une formule sont des nombres \emph{exacts}.
\begin{enumerate}[label=\alph*)]
  \item Calcule $M(\ce{H2O})$ avec le bon nombre de décimales.
  \item Déduis le nombre de moles dans $m=\qty{9,00}{\gram}$ d'eau, avec le bon nombre de CS.
\end{enumerate}
\end{exo}

\begin{exo}[quand un facteur est exact]
Un nombre exact (comptage, partage) a une infinité de CS et ne limite jamais le résultat.
\begin{enumerate}[label=\alph*)]
  \item On aligne $5$ flacons contenant chacun $\qty{2,50}{\gram}$ de sel. Masse totale ?
  \item On répartit $\qty{250,0}{\gram}$ de poudre en $4$ portions \emph{égales}. Masse d'une portion ?
\end{enumerate}
\end{exo}

\begin{exo}[pH : le bon nombre de décimales]
Utilise la règle «~décimales du pH $=$ CS de la concentration~».
\begin{enumerate}[label=\alph*)]
  \item $[\ce{H3O+}]=2,5\times10^{-4}\ \si{\mol\per\litre}$ : calcule le pH.
  \item $[\ce{H3O+}]=1,00\times10^{-3}\ \si{\mol\per\litre}$ : calcule le pH.
  \item $\mathrm{pH}=2,00$ : remonte à $[\ce{H3O+}]$.
\end{enumerate}
\end{exo}

\begin{exo}[applique la bonne règle]
Pour chaque calcul, donne le résultat correctement arrondi.
\begin{enumerate}[label=\alph*)]
  \item $45,2 + 3,08$
  \item $6,4 \times 0,125$
  \item $78,5 \div 2,50$
  \item $12,07 - 0,3$
\end{enumerate}
\end{exo}

\end{document}
